\section{Related Work}
\label{sec:related_work}

Our research sits at the intersection of weight-space model merging, dynamic Mixture-of-Experts routing, and parameter-efficient multi-task adaptation. 

\subsection{Parameter-Space Model Merging}
Directly combining multiple specialized networks in weight space has become an active area of study to facilitate multi-task adaptation without retraining~\cite{he2016deep, modelsoups}. Early techniques focused on linear weight averaging or interpolation, such as Model Soups~\cite{wortsman2022model} and Federated Learning aggregation~\cite{federated_learning}. To resolve parameters conflicts, recent SOTA techniques introduce structural fusions: Task Arithmetic~\cite{ilharco2022editing} subtracts the pre-trained base model weights to construct "task vectors" which are then scaled and added; Ties-Merging~\cite{yadav2023ties} addresses sign disagreements and parameter interference by resolving sign conflicts and keeping only dominant weights; DARE~\cite{dare} randomly prunes and rescales weights before merging. However, these methods calculate a static set of coefficients that remain identical across all samples, leading to performance collapse when experts represent highly divergent or conflicting task domains.

\subsection{Dynamic Routing and Mixture-of-Experts}
Our work is also conceptually related to Mixture-of-Experts (MoE) architectures~\cite{moe}, which route token-level activations to specialized sub-networks or "experts" dynamically. MoEs rely on routing layers trained via supervised joint-training or large-scale pre-training~\cite{switch_transformers, jacobs1991adaptive}. However, traditional MoEs require specialized, routing-compatible architectures and carry immense parameter and memory overhead during inference since all expert parameters must remain resident. In contrast, dynamic model merging performs on-the-fly fusion of expert weights \emph{before} executing the forward pass, combining multiple expert models into a single standard model instance during inference, thereby carrying zero additional memory or computational latency overhead in the backbone layers.

\subsection{Quantum vs. Classical Dynamic Merging}
Recently, dynamic model merging has evolved to address the rigid nature of static fusions. AdaMerging~\cite{adamerging} adaptively tunes merging coefficients on an offline calibration set, though the coefficients remain static at inference time. To realize sample-dependent merging, Quantum Wavefunction Superposition Merging (QWS-Merge)~\cite{qwsmerge} routes parameters using quantum-inspired phase superposition gating. More recently, BSigmoid-Router established a classical alternative using independent, bounded sigmoidal activations, which eliminates the competitive Softmax constraints. 

Despite these advances, existing dynamic merging methods are limited by their routing input feature design. They universally apply global average pooling to the token features of the backbone to obtain a flat representation vector before mapping it to task logits. This average pooling operation is a critical empirical flaw: it collapses spatial resolution, making the routing decisions vulnerable to spatial occlusions, noise, and patch corruptions. Furthermore, during multi-task batch processing, average pooling across heterogeneous samples averages out the task-specific feature signatures, causing "heterogeneity collapse" where the dynamic coefficients default to flat uniform fusions. CAM-Router resolves this fundamental gap by preserving the un-pooled token sequence and using a cross-attention mechanism with trainable task queries.
